We conducted three rounds of adversarial red-teaming on the detection
methodology before running the clean experiment. We report every critical
finding and its resolution as part of the scientific record. This section
is not an afterthought---it is the most important methodological
contribution of this paper.

\subsection{Round 1: Fundamental Design Flaws (Prior to Clean Experiment)}

The original Oracle Loop detection experiment (v1) used instructed
conditions: the system prompt told the model to ``answer honestly'' or
``make up a plausible-sounding answer.'' Red-teaming identified four
critical issues:

\begin{redteam}
\textbf{C1: SVD features are nonlinear in token count.} The relationship
between sequence length and singular value structure is nonlinear
(eigenvalues scale approximately as $\mathcal{O}(n^{1/2})$ to
$\mathcal{O}(n)$ depending on the spectral regime). Linear FWL
residualization is insufficient. \textit{Resolution:} Marchenko-Pastur
correction replaces post-hoc regression with a dimension-aware theoretical
threshold (\cref{eq:mp_threshold}).
\end{redteam}

\begin{redteam}
\textbf{C2: Behavioral classifier has asymmetric noise.} Pattern-matching
classifiers for hedging vs.\ confabulation have different false-positive
rates per category. A model that says ``I believe the answer is X'' may be
classified as HEDGED or CONFABULATED depending on arbitrary pattern
thresholds. \textit{Resolution:} Added AMBIGUOUS category for epistemic
hedges; excluded from primary analysis.
\end{redteam}

\begin{redteam}
\textbf{C3: Encoding features leak into generation analysis.} When the
full KV cache is used for feature extraction, encoding-phase structure
dominates generation-phase signal. Instructed conditions create different
system prompts, making encoding features a trivial classifier.
\textit{Resolution:} Generation-only cache slicing with windowed
extraction. Single system prompt for all conditions.
\end{redteam}

\begin{redteam}
\textbf{C4: Non-random prompt-to-behavior assignment.} Instructed
conditions assign prompts to behaviors by construction. The classifier
may learn prompt features, not cognitive geometry.
\textit{Resolution:} Post-hoc behavioral classification from response
content. Same prompt pool for HEDGED and CONFABULATED.
\end{redteam}

\subsection{Round 2: Implementation Bugs in MP Correction}

After implementing the MP correction, a second red-team pass identified:

\begin{redteam}
\textbf{CRITICAL: Row permutation is a no-op for SVD.} The empirical null
function (\texttt{compute\_empirical\_null}) originally permuted
\emph{rows} of the flattened key matrix. But for any permutation matrix
$P$, $(PX)^T(PX) = X^T P^T P X = X^TX$---row permutation preserves the
Gram matrix exactly. The function returned the top eigenvalue of the
original data as the ``null,'' making all MP features collapse to
near-constants and all invariance diagnostics trivially pass.
\textit{Resolution:} Column-wise permutation (shuffle each column
independently via \texttt{np.argsort(np.random.random((n,p)), axis=0)}).
This destroys covariance structure while preserving marginal distributions,
the standard approach for spectral null testing \citep{buja2009comment}.
\end{redteam}

\begin{redteam}
\textbf{Regression diagnostic used wrong regressors.} The invariance
diagnostic regressed MP features on both \texttt{gen\_window} (capped at
80) and $\log(n_{\text{generated}})$. Since \texttt{gen\_window} is
near-constant for full-window trials, the two regressors were collinear.
\textit{Resolution:} Restrict to full-window trials, single
$\log(n_{\text{generated}})$ regressor.
\end{redteam}

\begin{redteam}
\textbf{FWL belt used wrong confound variable.} The belt-and-suspenders
FWL on MP features used \texttt{n\_generated\_tokens} as the confound, but
MP features are computed on a windowed cache capped at 80 tokens.
The correct confound is \texttt{gen\_window} (actual window size used).
\textit{Resolution:} Updated to use \texttt{gen\_window} as confound for
MP features; \texttt{n\_generated\_tokens} for raw features.
\end{redteam}

\begin{redteam}
\textbf{mp\_norm\_per\_token normalization error.} Division by
\texttt{seq\_len} instead of \texttt{n\_heads $\times$ seq\_len},
underestimating the denominator by a factor of $h$.
\textit{Resolution:} Fixed to divide by $n = h \times T$.
\end{redteam}

\subsection{Round 3: Validation}

A final red-team pass verified all fixes:

\begin{enumerate}[nosep]
  \item Column-wise permutation produces genuinely different SVD spectra.
  \item Empirical $\lambda_+$ differs from theoretical by orders of
    magnitude (confirming serial correlation in KV-cache keys).
  \item Single-regressor diagnostic on full-window trials: 4/5 features
    $R^2 < 0.03$.
  \item FWL confound variables match feature computation scope.
  \item All 4-tuple format entries in the analysis pipeline parse
    correctly.
\end{enumerate}

\subsection{Prior Scaler Leak (v1)}

Our original Oracle Loop experiment (v1) contained a subtler bug:
the \texttt{StandardScaler} was fitted on all data before the LOO loop,
then the same scaled features were used for train and test in each fold.
This leaks global mean and variance into every test prediction.

We recomputed the v1 results with correct scaler-in-LOO
(\cref{tab:scaler_leak}):

\begin{table}[H]
\centering
\caption{Impact of scaler leak on v1 results. $\Delta$ shows the
  correction when the scaler is fitted inside LOO.}
\label{tab:scaler_leak}
\begin{tabular}{lrrr}
\toprule
Comparison & Leaked LOO & Corrected LOO & $\Delta$ \\
\midrule
honest vs confab       & 0.934 & 0.930 & $-0.004$ \\
honest vs deceptive    & 0.904 & 0.899 & $-0.005$ \\
confab vs deceptive    & 0.846 & 0.841 & $-0.005$ \\
\bottomrule
\end{tabular}
\end{table}

The leak's impact was small ($< 0.01$ AUROC) because LOO removes only one
sample, so $n-1$ samples closely approximate the full-data scaler. We
report this for completeness and to establish that our v1 detection claims
survive the correction.

\subsection{Lessons}

Three lessons from the red-team process:

\begin{enumerate}[nosep]
  \item \textbf{Mathematical invariance proofs matter.} The row-permutation
    bug would have made every MP feature collapse silently. No amount of
    ``the numbers look reasonable'' checking would have caught it---only
    understanding the algebraic structure of SVD revealed the problem.
  \item \textbf{Defensive design beats defensive testing.} The MP
    correction eliminates the token-count confound by construction rather
    than correcting for it post-hoc. This is inherently more robust than
    any regression-based fix.
  \item \textbf{Report the bugs.} The temptation to quietly fix errors and
    present only the corrected results is strong. We resist it. The bugs
    and their fixes are the methodology.
\end{enumerate}
